Three points follow from these results. First, multilingual variation is not the same as human validity. A model can look plausible under multilingual prompting yet still diverge substantially from the Japanese human reference when calibration is evaluated directly. This is the main reason we treat Japanese human alignment as the primary endpoint and EN/ZH behavior as robustness analysis rather than as co-equal human-grounded evaluation.

Second, a compact benchmark can still produce operationally useful screening information. The two-axis view in Figure~\ref{fig:tradeoff} is especially helpful here. If an application mainly requires primary-language closeness to Japanese human readers, Claude Sonnet 4.5 is the strongest choice in this panel. If multilingual stability under translated conditions is the stronger requirement, Qwen3.6 Plus and GPT-5.4 are the most attractive. The benchmark therefore supports model selection without pretending that there is one universally best model for every affective-language use case.

\begin{table}[t]
\caption{Practical screening interpretation of the six-model panel.}
\label{tab:screening}
\centering
\scriptsize
\begin{tabular}{@{}p{0.24\columnwidth}p{0.30\columnwidth}p{0.29\columnwidth}@{}}
\toprule
Model & Strength & Caution \\
\midrule
Claude Sonnet 4.5 & Best Japanese MAE & Only mid-level average multilingual drift \\
DeepSeek V3.2 & Second-best Japanese MAE & Largest average drift; English especially unstable \\
Qwen3.6 Plus & Lowest average EN/ZH drift & Japanese alignment below the top two models \\
GPT-5.4 & Highest Japanese Pearson and Spearman; near-best average drift & Not the lowest Japanese MAE \\
Grok 4.20 & Mid-range on both axes & No leading metric in the six-model panel \\
Gemini 2.5 Pro & Moderate average drift & Largest Japanese MAE in the panel \\
\bottomrule
\end{tabular}
\end{table}

Table~\ref{tab:screening} summarizes the same point in operational terms. The benchmark does not identify a single dominant model; instead, it exposes which trade-off each model is making in a compact and reviewable form. That makes the protocol useful for early-stage model narrowing before a larger validation campaign is launched.

Third, the paper clarifies how this ICECCME contribution differs from our earlier LLM-only framework paper. The previous study was designed to characterize emotional individuality under persona-assigned temperatures across a broad 36-model landscape \cite{shirahama2026mdpi}. The present paper instead fixes a neutral-reader primary condition and uses human alignment as the primary endpoint. This shift makes the paper more appropriate for a compact engineering conference contribution focused on benchmark design and cross-language robustness rather than on multidimensional individuality analysis.

This neutral-reader constraint is methodological, not merely stylistic. Once role-playing or creativity prompts are introduced, it becomes harder to know whether score changes reflect better human alignment, a different interpretive stance, or prompt-induced variance. For a human-grounded pilot benchmark, anchoring the main endpoint to one neutral condition makes the comparison easier to interpret and reproduce across providers. Richer persona-conditioned evaluation remains worthwhile, but it should be layered onto a stable human-reference baseline rather than replacing it.

The benchmark should also be understood as a stage-1 filter rather than as a deployment certificate. A model that scores well here is better treated as a candidate for deeper follow-up than as automatically suitable for education, counseling, or literary-analysis support. In those higher-stakes settings, one would still want broader text coverage, more diverse readers, richer explanation analysis, and application-specific human review. The present contribution is deliberately modest: it shows that even a small human-grounded protocol can distinguish calibration quality from multilingual stability in a way that purely model-versus-model comparison cannot.

The main limitations are equally clear. The benchmark uses only three texts, one Japanese human reference condition, and one neutral-reader primary run. English and Chinese were introduced as robustness conditions, not as human-grounded evaluation conditions in their own right. In addition, provider-compatibility adjustments were required at the transport/schema layer during data collection, although the task definition, source texts, and scoring dimensions were kept unchanged. The six-model panel is intentionally contemporary and compact rather than exhaustive. These limitations should be read as scope conditions rather than as weaknesses of the core human-grounded comparison.

A natural next step is therefore not to widen the current claims, but to layer new evidence in a controlled order: more texts, larger human reference samples, and matched human references for English and Chinese. Those extensions would allow future work to ask whether the present Japanese-centered screening picture remains stable once multilingual human calibration is available.
